% Corrected Corollary 3.1.3: invertible ideals at conductor primes.

\begin{corollary}[Corrected Corollary 3.1.3]\label{cor-3-1-3}
\lean{QuadraticOrder.invertibleIdealCount_prime_pow_trivial,
  QuadraticOrder.invertibleIdealCount_prime_pow_ramified,
  QuadraticOrder.invertibleIdealCount_prime_pow_inert,
  QuadraticOrder.invertibleIdealCount_prime_pow_split}
\uses{notation, lem-3-2-4, lem-3-2-7, lem-3-2-8,
  infra-index, infra-invertibility}
Let $d=Df^2$, where $D<0$ is fundamental, and let the prime $p$ divide $f$.
Put $w=v_p(f)$ and put $v=v_p(d)$ for odd $p$, while
$v=v_2(d)-2=v_2(d/4)$ for $p=2$.  The number $C_d^\times(p^m)$ of invertible
ideals of $\mathcal O_d$ of index $p^m$ is as follows.

If $m<v$, or if $m=v=2w-2$, then
\[
C_d^\times(p^m)=
\begin{cases}p^{m/2},&m\text{ even},\\0,&m\text{ odd}.
\end{cases}
\]
If $v<m$, or if $m=v$ and $2w\leq v$, then
\[
C_d^\times(p^m)=
\begin{cases}
p^w,&(D/p)=0,\\
p^w+p^{w-1},&(D/p)=-1\text{ and }m\text{ even},\\
0,&(D/p)=-1\text{ and }m\text{ odd},\\
(p^w-p^{w-1})(m+1-2w),&(D/p)=1.
\end{cases}
\]
In particular the boundary $m=v=2w-2$, which occurs for $p=2$ and odd $D$,
belongs to the first regime (ERRATA E5).
\end{corollary}

\begin{proof}
\uses{lem-3-2-4, lem-3-2-7, lem-3-2-8,
  infra-index, infra-invertibility}
Write
$g(X)=X^2-dX+(d^2-d)/4$.  The normal-form bijection writes every ideal of
index $p^m$ uniquely as
\[
\mathfrak a_{k,A}=p^k\mathbb Z\oplus
p^{m-k}(\tau-A)\mathbb Z,qquad
\lceil m/2\rceil\leq k\leq m,
\]
where, for $s=2k-m$, the parameter $A$ is a root of $g$ modulo $p^s$.
Put $r=m-k$.  Lemma~\ref{lem-3-2-7} says that the contraction of the local
principal ideal generated by $p^r(\tau-A)$ has parameters
$r+\nu,2r+\nu$, with $\nu=v_p(g(A))$.  Uniqueness and
Lemma~\ref{lem-3-2-8} therefore give, for $s>0$,
\[
\mathfrak a_{k,A}\text{ invertible}\quad\Longleftrightarrow\quad
v_p(g(A))=s.
\tag{1}
\]
For $s=0$ the ideal is the principal ideal $p^{m/2}\mathcal O_d$.

Let $E_s$ count the exact roots in (1), and set $E_0=1$.  Thus
\[
C_d^\times(p^m)=
\sum_{\substack{0\leq s\leq m\\s\equiv m\pmod2}}E_s.\tag{2}
\]
Completing the square translates $g$ to $x^2-q$, where $v_p(q)=v$ (for
$p=2$, $q=d/4$ is integral).  If $0<s<v$, exact valuation forces
$s=2t$ and $x=p^ty$ with $y$ a unit; hence
\[
E_{2t}=p^t-p^{t-1},\qquad E_{2t+1}=0.
\tag{3}
\]
The even terms telescope to
$1+\sum_{t=1}^h(p^t-p^{t-1})=p^h$, proving the first regime.

Write $q=p^vu$.  At and above the boundary, reducing $y^2-u$ after extracting
$p^{\lfloor v/2\rfloor}$ gives the following exhaustive alternatives.
If $v=2w+1$, the boundary contributes $p^w$ and an odd-valuation unit cannot
be a square above it.  If $v=2w$ and $p$ is odd, the sum through the boundary
is $p^w+p^{w-1}$ for nonsquare $u$ and $p^w-p^{w-1}$ for square $u$; the
former has no tail, while every permitted tail exponent in the latter case
contributes $2(p^w-p^{w-1})$.  If $p=2$ and $v=2w$ with
$u\equiv3\pmod4$, the even boundary and the possible odd boundary each give
$2^w$, and there is no later root.  Finally, if $p=2$, $D$ is odd, and
$v=2w-2$, the even sum through $v$ is $2^{w-1}$; for
$u\equiv5\pmod8$ only exponent $v+2$ contributes (with $2^w$ classes), while
for $u\equiv1\pmod8$ the first two tail exponents contribute zero and every
permitted exponent from $v+3$ onward contributes $2^w$.

These counts use only that odd squares are $1$ modulo $8$ and that, for odd
$p$, a square unit has two simple roots while a nonsquare has none.  Inserting
them into (2) gives $p^w$ in the ramified case,
$p^w+p^{w-1}$ for even inert $m$ and zero for odd inert $m$, and
\[
(p^w-p^{w-1})(m+1-2w)
\]
in the split case.  The $v=2w-2$ equality remains in the telescoping block,
which is exactly the corrected boundary asserted above.
\end{proof}

\begin{remark}\label{cor-3-1-3-errata}
\uses{cor-3-1-3}
The thesis places every equality $m=v$ in the second regime.  At
$d=-112$, $w=2$, $m=v=2$, this makes its split expression negative; the
correct count is $2$.  The regression rows for $d=-48,-112,-192$ enforce the
corrected boundary.
\end{remark}
